\documentclass[../../../../main.tex]{subfiles}
\begin{document}

\begin{flushright}
\footnotesize\textit{【自動文字起こし・要確認】}
\end{flushright}

{\bf [解]}
全ての操作の後，$N+1$に赤が入っているのは，$N$回目の操作後に赤の入っている箱を$T$として，$N+1$回目で$T$と交換する時である（$T\ne N+1$）．$T\ne N+1$となる確率は，排反を考えて，
\begin{align*}
1-\left(\frac{N-1}N\right)^N
\end{align*}

だからもとめるのは
\begin{align*}
\frac1N\left\{1-\left(\frac{N-1}N\right)^N\right\}．
\end{align*}

\newpage

\end{document}
